\clearpage
\scanpage{3}

A) \uline{Descente et passage au quotient}

\S\,1\quad Descente fidèlement plate

\S\,2\quad \deleted{Passage au quoti\unsure{ent|.} \unsure{Généralités} sur les foncteurs sur} \inserted{Topologies sur $(\mathrm{Sch})_{/S}$. Propriétés générales des foncteurs sur $(\mathrm{Sch})_{/S}$.} \margin{[Intervertir \uline{1} et \uline{2} ??]}
\par\noindent\hfill\deleted{$\unsure{(\mathrm{Sch})_{/S}|(\mathrm{Ens})_{/S}}$}

\S\,3\quad Passage au quotient plat.

\S\,4\quad Descente non plate\quad [Application \unsure{à|:} \unsure{Ram.|Rau.|Rem.} \unsure{Samuel|Serreil} \unsure{généralisé} \unsure{\ldots|---}

\inserted{Appl. à platitude d'une \unsure{adh.} \unsure{schématique} et \unsure{lemme} de \deleted{\unsure{H}}\unsure{Hironaka}.}

\S\,5\quad Passage au quotient non plat. [\unsure{Passage} \unsure{aux|au} \unsure{quot.|quotients} \unsure{cat.|catég.} \unsure{gén}]

\S\,6\quad Descente \unsure{de|des} schémas relatifs étales, ----

[\unsure{Év\textsuperscript{t}} descente des \deleted{\unsure{morphismes}}\inserted{faisceaux} étales ?? \ldots]

B) \deleted{\uline{Propriétés générales des foncteurs sur} $(\mathrm{Sch})_{/S}$.}

\deleted{\unsure{\S\,7}}\quad \uline{Méthodes \inserted{de construction} non projectives}

\deleted{\S\,7\quad Propriétés générales \unsure{de} \unsure{foncteurs}, et propriétés de permanence}\quad \deleted{\inserted{\unsure{diverses}}}

\S\,7\quad \uline{Proreprésentabilité}\unsure{:|.} (et variantes).

\tikz[baseline=(p.base),remember picture]{\node[inner sep=1pt] (p) {\S\,\deleted{\unsure{8}}\unsure{9}};}\quad
\uline{Invariants différentiels} \deleted{\unsure{de|du} \unsure{div.}} \uline{\unsure{tangts|tang.} et normaux}

\tikz[baseline=(q.base)]{\node[inner sep=1pt] (q) {\S\,10};\draw (q.center) ellipse (1.25em and .9em);}\quad
\margin{\fbox{Après \S\,12 ??}}Propriétés générales \unsure{de|des} foncteurs, \uline{propriétés de permanence}
\margin{$\left\{\begin{gathered}\text{Foncteurs \uline{\unsure{excellent|excellents}} sur S}\\ \text{\unsure{morphismes|impliquent} excellent}\\ \text{??}\end{gathered}\right.$}

\tikz[baseline=(q.base)]{\node[inner sep=1pt] (q) {\S\,\deleted{\unsure{10|11}}11};\draw (q.center) ellipse (1.65em and .9em);}\quad
% slanted
\margin{\fbox{\parbox{.90\marginparwidth}{\raggedright\unsure{Commun au|Convient au} \unsure{par.} \unsure{20|éd} ? Cas Greenberg ?}}}
Représentabilités ``élémentaires'' \unsure{des|du} type $\prod_{\deleted{\unsure{X}}Y/S}X/Y$.

\inserted{cf SGAD $\underline{\overline{\mathrm{VIII}}}$ 6 et $\underline{\overline{\mathrm{XI}}}$ 6.8.}

\tikz[baseline=(r.base),remember picture]{\node[draw,rounded corners,inner sep=1pt] (r) {\S\,\unsure{8}\deleted{\unsure{1|11}}};}
\quad Construction de Modules cohérents
\figure{\begin{tikzpicture}[remember picture,overlay]\draw[->] (r.north west) to[out=125,in=190] (p.north);\end{tikzpicture}}

\S\,12\quad \deleted{\unsure{Représentabilité}}\inserted{Foncteurs} \unsure{repr.} par morphismes \inserted{\unsure{séparés|xprés}} loc. \unsure{pres.} \unsure{fi|finie}

\S\,13\quad Applications : des \unsure{th\textsuperscript{es}|th\textsuperscript{s}} \deleted{de} représentabilité relatifs. (y compris \unsure{(}pour les \unsure{esp.|es} foncteurs étales\ldots\unsure{,|)}

% archive sheet number: 2
